\documentclass[12pt,a4paper]{article}

\usepackage[utf8]{inputenc}
\usepackage[T1]{fontenc}
\usepackage[italian]{babel}
\usepackage{geometry}
\usepackage{array}
\usepackage{longtable}
\usepackage{booktabs}
\usepackage{graphicx}
\usepackage{hyperref}
\usepackage{enumitem}
\usepackage{xcolor}
\usepackage{float}
\usepackage{tabularx}
\usepackage{multirow}
\usepackage{pgfplots}
\usepackage{ragged2e}

\pgfplotsset{compat=1.18}

\renewcommand{\arraystretch}{1.3}

\setlength{\parindent}{0pt}

\geometry{margin=2.5cm}

\title{\textbf{Milestone 3 Report}\\
TrentoParking}
\author{Team 12}
\date{17 Maggio 2026}

\begin{document}

\maketitle
\tableofcontents
\newpage

%-------------------------------------------------
\section{Sezione Introduttiva}

\subsection{Team Members}

\begin{table}[H]
\centering
\begin{tabular}{|p{6cm}|p{2.5cm}|p{7cm}|}
\hline
\textbf{Nome e Cognome} & \textbf{Matricola} & \textbf{Account Github} \\
\hline
David Dorobantu & 234467 & https://github.com/DavidDorobantu \\
\hline
Riccardo Gonzato & 246476 & https://github.com/Gonzo04 \\
\hline
Matteo Sepa & 243283 & https://github.com/sepamatteo \\
\hline
\end{tabular}
\caption{Membri del team}
\end{table}

\subsection{Project Idea}

TrentoParking è un'applicazione web dedicata al noleggio di posti auto privati.
L'obiettivo del progetto è permettere ai proprietari di parcheggi privati di mettere a disposizione il proprio posto auto quando non utilizzato, mentre gli utenti possono cercare e prenotare rapidamente un parcheggio disponibile tramite una mappa interattiva.
L'applicazione vuole ridurre il tempo necessario alla ricerca di parcheggio, migliorare la mobilità urbana e ottimizzare l'utilizzo degli spazi privati inutilizzati.

\subsection{Links}

\begin{itemize}
    \item Repository Github: \url{https://github.com/Gonzo04/TrentoParking}
    \item Link Apiary: \url{https://trentoparking.docs.apiary.io/#}
\end{itemize}

%-------------------------------------------------
\section{Sezione Generale}

\subsection{Strategia di Branching}

Il team ha adottato una strategia di branching basata sul modello \textit{GitFlow semplificato senza develop}.

I branch utilizzati sono:

\begin{itemize}
    \item \textbf{main}: contiene esclusivamente versioni stabili e funzionanti dell'applicazione.
    \item \textbf{feature/*}: branch dedicati allo sviluppo di specifiche funzionalità.
\end{itemize}

Ogni nuova funzionalità viene sviluppata in un branch separato e successivamente integrata in \textit{main} tramite pull request.
Prima del merge vengono effettuati:

\begin{itemize}
    \item revisione del codice;
    \item verifica della compilazione;
    \item controllo dei conflitti;
    \item test delle API e delle funzionalità implementate.
\end{itemize}

Questa strategia ha permesso al nostro team di lavorare in parallelo riducendo conflitti e mantenendo una maggiore stabilità del codice.

\subsection{Product Backlog}

\small

\begin{longtable}{|p{0.7cm}|p{2.5cm}|p{5.8cm}|p{1.3cm}|p{4.2cm}|}
\hline
\textbf{ID} & \textbf{Name} & \textbf{User story} & \textbf{Imp.} & \textbf{How to Demo} \\
\hline
1 & Registrazione &
Come utente voglio registrarmi all'applicazione tramite email e password così da poter utilizzare i servizi di TrentoParking. &
100 &
L'utente completa la registrazione e riceve l'email di verifica. \\
\hline

2 & Verifica Email &
Come utente voglio verificare il mio account tramite email così da poter accedere all'applicazione in sicurezza. &
60 &
L'utente clicca il link ricevuto via email e l'account viene attivato. \\
\hline

3 & Login &
Come utente voglio effettuare il login così da poter accedere alla piattaforma. &
100 &
L'utente inserisce le credenziali corrette e accede alla home. \\
\hline

4 & Password reset &
Come utente voglio recuperare la password dimenticata così da poter riottenere accesso al mio account. &
60 &
L'utente richiede il reset password e riceve l'email dedicata. \\
\hline

5 & Mappa interattiva &
Come utente voglio visualizzare una mappa interattiva così da poter trovare parcheggi disponibili. &
60 &
Dopo il login viene mostrata correttamente la mappa principale. \\
\hline

6 & Visualizzazione Parcheggio &
Come utente voglio vedere i posti auto disponibili sulla mappa così da scegliere il migliore. &
80 &
I marker dei parcheggi vengono caricati dinamicamente sulla mappa. \\
\hline

7 & Filtro per raggio &
Come utente voglio filtrare i parcheggi disponibili in base a un raggio di distanza così da visualizzare solo quelli vicini alla mia posizione. &
30 &
L'utente seleziona un raggio di ricerca e la mappa mostra solamente i parcheggi inclusi nell'area scelta. \\
\hline

8 & Menu parcheggi card &
Come utente voglio visualizzare i parcheggi disponibili anche in una lista laterale così da confrontare rapidamente le opzioni disponibili. &
80 &
Accanto alla mappa viene mostrata una sidebar con la lista dei parcheggi disponibili sincronizzata con i marker presenti sulla mappa. \\
\hline

9 & Selezione parcheggio dalla mappa &
Come utente voglio selezionare un parcheggio dalla mappa o dalla lista laterale. &
50 &
L'utente seleziona un parcheggio e visualizza i dettagli. \\
\hline

10 & Prenotazione &
Come utente voglio prenotare un posto auto disponibile. &
100 &
L'utente conferma la prenotazione tramite l'applicazione. \\
\hline

11 & Recensioni &
Come utente voglio lasciare recensioni ai proprietari dei parcheggi. &
40 &
L'utente inserisce una recensione dopo la fine della prenotazione. \\
\hline

12 & Messaggi &
Come utente voglio comunicare con il proprietario del parcheggio. &
40 &
L'utente apre una chat e invia messaggi al proprietario. \\
\hline

\end{longtable}

\normalsize
\vspace{1cm}

\subsection{Definizione di Done}

Una funzionalità viene considerata \textit{Done} quando soddisfa tutti i seguenti criteri:

\begin{itemize}
    \item il codice è stato implementato e caricato su Github;
    \item il codice compila senza errori;
    \item sono stati eseguiti test manuali e/o automatici;
    \item il codice è stato revisionato da almeno un altro membro del team;
    \item non sono presenti bug critici noti;
    \item la user story associata soddisfa i requisiti definiti durante lo sprint planning.
\end{itemize}

%-------------------------------------------------
\section{Sprint 1}

\subsection{Goal}

L'obiettivo principale dello Sprint 1 è stato quello di progettare e sviluppare l'infrastruttura iniziale dell'applicazione TrentoParking.

Durante questo sprint il team si è concentrato principalmente su:

\begin{itemize}
    \item configurazione dei repository Github;
    \item definizione dell'architettura frontend e backend;
    \item implementazione del sistema di autenticazione;
    \item registrazione utenti con verifica email;
    \item implementazione del login;
    \item implementazione del reset password;
    \item creazione della schermata principale con mappa interattiva;
    \item progettazione preliminare del sistema di prenotazione.
\end{itemize}

Lo sprint aveva inoltre lo scopo di creare una base solida e scalabile per lo sprint successivo.

\subsection{Sprint Planning}

Durante lo Sprint Planning il team ha analizzato le user story prioritarie e stimato il carico di lavoro.

\small

\renewcommand{\arraystretch}{1.25}

\setlength{\LTleft}{0pt}
\setlength{\LTright}{0pt}

\begin{longtable}{
|>{\centering\arraybackslash}m{1.8cm}
|>{\centering\arraybackslash}m{1.7cm}
|>{\raggedright\arraybackslash}m{4.7cm}
|>{\raggedright\arraybackslash}m{3.8cm}
|>{\centering\arraybackslash}m{1.4cm}
|>{\centering\arraybackslash}m{0.7cm}|}
\hline

\multicolumn{3}{|c|}{} &
\multicolumn{3}{c|}{\textbf{Sprint Backlog (Sprint Planning)}} \\
\cline{4-6}

\multicolumn{3}{|c|}{} &
\textbf{Task} & \textbf{Volunt.} & \textbf{Est.} \\
\hline
\endfirsthead

\hline
\multicolumn{3}{|c|}{} &
\multicolumn{3}{c|}{\textbf{Sprint Backlog (continued)}} \\
\cline{4-6}

\multicolumn{3}{|c|}{} &
\textbf{Task} & \textbf{Volunt.} & \textbf{Est.} \\
\hline
\endhead

\multirow{30}{*}{\textbf{Sprint \#1}}

& \multirow{6}{*}{Sign up}
& \multirow{6}{4.8cm}{Come utente voglio registrarmi all'applicazione tramite email e password così da poter utilizzare i servizi di TrentoParking.}

& Design API registrazione & Riccardo & 2 \\
\cline{4-6}

&&
& Implementazione endpoint REST & Riccardo & 4 \\
\cline{4-6}

&&
& Progettazione database utenti & Riccardo & 2 \\
\cline{4-6}

&&
& Sviluppo UI registrazione & Riccardo & 4 \\
\cline{4-6}

&&
& Testing registrazione & Riccardo & 2 \\
\cline{4-6}

&&
& Deploy autenticazione & Riccardo & 1 \\
\cline{2-6}

& \multirow{5}{*}{Verifica \\ email}
& \multirow{5}{4.8cm}{Come utente voglio verificare il mio account tramite email così da poter accedere all'applicazione in sicurezza.}

& Configurazione servizio email & David & 3 \\
\cline{4-6}

&&
& Generazione token verifica & David & 2 \\
\cline{4-6}

&&
& Validazione token verifica & David & 2 \\
\cline{4-6}

&&
& Testing verifica email & David & 2 \\
\cline{4-6}

&&
& Deploy servizio email & David & 1 \\
\cline{2-6}

& \multirow{6}{*}{Login}
& \multirow{6}{4.8cm}{Come utente voglio effettuare il login così da poter accedere alla piattaforma.}

& Design API login & Riccardo & 2 \\
\cline{4-6}

&&
& Implementazione autenticazione JWT & Riccardo & 4 \\
\cline{4-6}

&&
& Sviluppo UI login & Riccardo & 3 \\
\cline{4-6}

&&
& Gestione errori autenticazione & Riccardo & 2 \\
\cline{4-6}

&&
& Testing login & Riccardo & 2 \\
\cline{4-6}

&&
& Deploy login service & Riccardo & 1 \\
\cline{2-6}

& \multirow{5}{*}{Password \\ reset}
& \multirow{5}{4.8cm}{Come utente voglio recuperare la password dimenticata così da poter riottenere accesso al mio account.}

& Generazione reset token & David & 2 \\
\cline{4-6}

&&
& Implementazione endpoint reset & David & 3 \\
\cline{4-6}

&&
& Sviluppo UI reset password & David & 3 \\
\cline{4-6}

&&
& Testing reset password & David & 2 \\
\cline{4-6}

&&
& Deploy reset password & David & 1 \\
\cline{2-6}

\pagebreak

& \multirow{8}{*}{Mappa \\ interattiva}
& \multirow{8}{4.8cm}{Come utente voglio visualizzare una mappa interattiva così da poter trovare parcheggi disponibili.}

& Studio integrazione mappe & Matteo & 2 \\
\cline{4-6}

&&
& Configurazione map provider & Matteo & 3 \\
\cline{4-6}

&&
& Sviluppo UI mappa & Matteo & 4 \\
\cline{4-6}

&&
& Gestione marker parcheggi & Matteo & 3 \\
\cline{4-6}

&&
& Implementazione filtro raggio & Matteo & 3 \\
\cline{4-6}

&&
& Sidebar lista parcheggi & Matteo & 4 \\
\cline{4-6}

&&
& Testing visualizzazione mappa & Matteo & 2 \\
\cline{4-6}

&&
& Deploy schermata home & Matteo & 1 \\
\hline

& \multirow{8}{*}{Parking \\ reservation}
& \multirow{8}{4.8cm}{Come utente voglio prenotare un posto auto disponibile così da assicurarmi un parcheggio prima dell'arrivo.}

& Design API prenotazioni & Matteo & 2 \\
\cline{4-6}

&&
& Implementazione endpoint prenotazione & Matteo & 4 \\
\cline{4-6}

&&
& Gestione disponibilità parcheggi & Matteo & 3 \\
\cline{4-6}

&&
& Salvataggio prenotazioni database & Matteo & 3 \\
\cline{4-6}

&&
& Sviluppo UI prenotazione & Matteo & 4 \\
\cline{4-6}

&&
& Conferma prenotazione utente & Matteo & 2 \\
\cline{4-6}

&&
& Testing prenotazioni & Matteo & 2 \\
\cline{4-6}

&&
& Deploy servizio prenotazioni & Matteo & 1 \\
\cline{2-6}

\multicolumn{5}{|r|}{\textbf{Total}} & \textbf{93} \\
\hline

\end{longtable}

\normalsize

\subsubsection{Burndown Chart}

Il Burndown Chart mostra l'andamento dello sprint e il completamento progressivo delle attività.

\begin{figure}[H]
\centering

\begin{tikzpicture}
\begin{axis}[
    width=15cm,
    height=8cm,
    xlabel={Giorni Sprint},
    ylabel={Ore di lavoro},
    xmin=1, xmax=21,
    ymin=0, ymax=95,
    xtick={1,3,5,7,9,11,13,15,17,19,21},
    ytick={0,10,20,30,40,50,60,70,80,90},
    legend pos=north east,
    grid=major
]

% Ideal Burndown
\addplot[
    mark=*,
]
coordinates {
    (1,93)
    (2,88)
    (3,84)
    (4,79)
    (5,75)
    (6,70)
    (7,66)
    (8,61)
    (9,57)
    (10,53)
    (11,48)
    (12,44)
    (13,39)
    (14,35)
    (15,30)
    (16,26)
    (17,22)
    (18,17)
    (19,13)
    (20,8)
    (21,0)
};

\addlegendentry{Ideale}

% Actual Burndown
\addplot[
    mark=square*,
]
coordinates {
    (1,93)
    (2,91)
    (3,88)
    (4,84)
    (5,80)
    (6,75)
    (7,70)
    (8,66)
    (9,60)
    (10,56)
    (11,50)
    (12,46)
    (13,42)
    (14,36)
    (15,31)
    (16,26)
    (17,21)
    (18,16)
    (19,10)
    (20,5)
    (21,0)
};

\addlegendentry{Vero}

\end{axis}
\end{tikzpicture}

\caption{Burndown Chart Sprint \1}

\end{figure}

\subsection{Test cases}

Di seguito sono riportati  20 test case relativi alle funzionalità dello Sprint \1 e al design delle funzionalità principali non ancora introdotte

\small
\begin{longtable}{|p{2.2cm}|p{0.6cm}|p{3.2cm}|p{2.8cm}|p{1.8cm}|p{3.2cm}|}
\hline
\textbf{User Story} & \textbf{ID} & \textbf{Description} & \textbf{Test Case Data} & \textbf{Precond.} & \textbf{Expected Result} \\ \hline
\endfirsthead
\hline
\textbf{User Story} & \textbf{ID} & \textbf{Description} & \textbf{Test Case Data} & \textbf{Precond.} & \textbf{Expected Result} \\ \hline
\endhead

% --- REGISTRAZIONE ---
\multirow{5}{=}{Registrazione} & 01 & Registrazione valida & email: user@test.it, pwd: Pwd123! & Nessuna & Account creato e redirect. \\ \cline{2-6}
 & 02 & Formato email errato & email: user@test, pwd: Pwd123! & Nessuna & Errore formato email. \\ \cline{2-6}
 & 03 & Email già esistente & email: existing@test.it, pwd: Pwd123! & Email in DB & Errore account duplicato. \\ \cline{2-6}
 & 04 & Campo email vuoto & email: [vuoto], pwd: Pwd123! & Nessuna & Messaggio campo richiesto. \\ \cline{2-6}
 & 05 & Password vuota & email: user@test.it, pwd: [vuoto] & Nessuna & Messaggio campo richiesto. \\ \hline

% --- LOGIN ---
\multirow{4}{=}{Login} & 06 & Login corretto & email: user@test.it, pwd: Pwd123! & Acc. attivo & Redirect alla Mappa. \\ \cline{2-6}
 & 07 & Password errata & email: user@test.it, pwd: wrong & Acc. attivo & Errore credenziali. \\ \cline{2-6}
 & 08 & Account non verificato & email: user@test.it, pwd: Pwd123! & Non verif. & Errore verifica email. \\ \cline{2-6}
 & 09 & Logout utente & Click pulsante Logout & Loggato & Sessione distrutta. \\ \hline

% --- RESET ---
\multirow{3}{=}{Password Reset} & 10 & Richiesta reset valida & email: user@test.it & Account esistente & Email di reset inviata. \\ \cline{2-6}
 & 11 & Reset con token valido & Nuova password: Pwd456! & Token OK & Password aggiornata in DB. \\ \cline{2-6}
 & 12 & Reset email inesistente & email: ignoto@test.it & Nessuna & Errore email non trovata. \\ \hline

% --- MAPPA & FILTRI ---
\multirow{4}{=}{Map \& Filter} & 13 & Caricamento mappa & Accesso Home page & Loggato & Mappa centrata su Trento. \\ \cline{2-6}
 & 14 & Visualizzazione marker & N/A & Dati in DB & Marker visibili in mappa. \\ \cline{2-6}
 & 15 & Filtro raggio 500m & Slider impostato su 0.5km & Marker OK & Filtraggio visivo marker. \\ \cline{2-6}
 & 16 & Nessun risultato & Raggio: 100m in zona vuota & Area vuota & Messaggio "Nessun risultato". \\ \hline

% --- SIDEBAR & PRENOTAZIONE ---
\multirow{2}{=}{Sidebar} & 17 & Sincronizzazione Lista & Filtro attivo su mappa & Marker OK & Lista aggiornata con mappa. \\ \cline{2-6}
 & 18 & Selezione da mappa & Click su marker specifico & Marker OK & Apertura schermata dettagli. \\ \hline
\multirow{2}{=}{Booking} & 19 & Prenotazione riuscita & Click "Prenota ora" & Posto Libero & Stato "Occupato" e conferma. \\ \cline{2-6}
 & 20 & Prenotazione fallita & Click su posto occupato & Occupato & Errore disponibilità. \\ \hline

\end{longtable}
\normalsize

A causa della dimensione della pagina non è possibilie fornire un livello di dettaglio sufficiente come indicato dal template, forniamo quindi un file .csv con il corretto livello di dettaglio:\\
\href{https://docs.google.com/spreadsheets/d/1MUQECu_8l71_o3VYYRUg6efop6524kugg39CWwuaUdM/edit?usp=sharing}{\textbf{Clicca qui per visualizzare il file dei Test Case}}.

\subsection{Sprint Review}

Durante la Sprint Review il team ha effettuato una dimostrazione delle funzionalità implementate.

Sono state presentate:

\begin{itemize}
    \item registrazione utenti;
    \item verifica email;
    \item login;
    \item reset password;
    \item navigazione verso la schermata principale;
    \item integrazione iniziale della mappa.
\end{itemize}

La discussione successiva si è concentrata principalmente su:

\begin{itemize}
    \item miglioramento dell'interfaccia utente;
    \item ottimizzazione delle API di autenticazione;
    \item gestione delle prenotazioni future;
\end{itemize}

È stato deciso che nel prossimo sprint verranno implementate le funzionalità principali relative alla ricerca e prenotazione dei parcheggi.

\subsection{Product Backlog Refinement}

Durante il refinement il team ha aggiornato il Product Backlog introducendo nuove user story emerse durante lo sviluppo.

In particolare sono state aggiunte:

\begin{itemize}
    \item visualizzazione dinamica dei parcheggi disponibili sulla mappa;
    \item lista laterale sincronizzata con la mappa;
    \item sistema di recensioni per i proprietari dei parcheggi;
    \item chat tra utente e proprietario;
    \item gestione disponibilità temporale dei parcheggi.
\end{itemize}

Inoltre alcune priorità sono state modificate per dare maggiore importanza alle funzionalità di prenotazione e visualizzazione dei posti auto.

\subsection{Sprint Retrospective}

Durante la retrospective il team ha analizzato gli aspetti positivi e negativi dello sprint.

\subsubsection{Aspetti positivi}

\begin{itemize}
    \item buona collaborazione tra i membri del team;
    \item utilizzo efficace di Github e delle pull request;
    \item comunicazione frequente durante lo sviluppo;
    \item completamento delle funzionalità di autenticazione;
    \item buona suddivisione dei task.
\end{itemize}

\subsubsection{Criticità riscontrate}

\begin{itemize}
    \item difficoltà iniziali nella configurazione dell'ambiente di sviluppo;
    \item sottostima del tempo necessario per alcune API;
\end{itemize}

\subsubsection{Miglioramenti per il prossimo sprint}

\begin{itemize}
    \item introdurre meeting tecnici più frequenti;
    \item migliorare la documentazione interna;
    \item completare il sistema di prenotazione;
    \item sviluppare recensioni e messaggistica;
    \item ottimizzare l'esperienza utente nella selezione dei parcheggi.
\end{itemize}

\end{document}
```
